% preamble.tex —— SY02-31《朗兰兹纲领 II》中文译注本 导言区
% 引擎：XeLaTeX（ctexbook）；要求 0 error / 0 undefined / 0 缺字
\documentclass[11pt,a4paper,UTF8,fontset=windows,openany]{ctexbook}

\usepackage{amsmath,amssymb,amsthm,mathtools}
\usepackage{mathrsfs}
\usepackage{stmaryrd}
\usepackage{bm}
\usepackage{geometry}
\geometry{a4paper,top=2.4cm,bottom=2.4cm,left=2.6cm,right=2.6cm,
  marginparwidth=1.35cm,marginparsep=0.18cm,headsep=0.5cm,footskip=0.9cm}
\usepackage{graphicx}
\usepackage{tikz}
\usetikzlibrary{cd,arrows.meta,positioning}
\usepackage{tikz-cd}
\usepackage[shortlabels]{enumitem}
\usepackage{booktabs,array,longtable,multirow}
\usepackage{xcolor}
\usepackage{fancyhdr}
\usepackage{caption}
\usepackage[colorlinks=true,linkcolor=blue!45!black,urlcolor=blue!45!black,
  citecolor=blue!45!black]{hyperref}
% 节标题里含数学时，PDF 书签串会出错；统一关掉书签，保留正文与目录超链接。
% （目录本身仍正常排版，因为那是排版结果而非 PDF 字符串。）
\hypersetup{bookmarks=false}

% ---------- 版式 ----------
\linespread{1.18}
\setlength{\parskip}{0.15em}
% 底本的编号是「篇.节.小节」三级（如 4.3.2），故须让 subsubsection 带编号。
\setcounter{secnumdepth}{3}
% 中文正文不靠连字符断行，长数学串最容易压出页边；给 TeX 三轮加宽容忍度。
\setlength{\emergencystretch}{3em}
\tolerance=3000
\hbadness=4000
\ctexset{
  chapter/format = \Large\bfseries\centering,
  chapter/name = {},
  chapter/number = {},
  chapter/beforeskip = 1.2em,
  chapter/afterskip = 1.0em,
  section/format = \large\bfseries\raggedright,
  subsection/format = \normalsize\bfseries\raggedright,
  subsubsection/format = \normalsize\bfseries\raggedright,
}
\pagestyle{fancy}
\fancyhf{}
\fancyhead[CE]{\small 朗兰兹纲领 II：朗兰兹对应的几何化}
\fancyhead[CO]{\small\leftmark}
\fancyfoot[C]{\small\thepage}
\renewcommand{\headrulewidth}{0.3pt}
\fancypagestyle{plain}{\fancyhf{}\fancyfoot[C]{\small\thepage}\renewcommand{\headrulewidth}{0pt}}

% 每篇（章）内小节号从 1 重新起，忠实原文各文独立编号
\renewcommand{\thesection}{\arabic{section}}
\renewcommand{\thesubsection}{\thesection.\arabic{subsection}}
\renewcommand{\thesubsubsection}{\thesubsection.\arabic{subsubsection}}

% ---------- 页锚：每页原文页码标在切口侧 ----------
\newcommand{\ysrc}[1]{\marginpar{\raggedleft\scriptsize\color{black!55}#1}%
  \hypertarget{src:#1}{}}

% ---------- 定理类环境：编号沿用原文，手工传入 ----------
% 用法：\begin{thm}{1.2.3} ... \label{thm:xxx} ... \end{thm}
% \ref{thm:xxx} 得到 1.2.3（已用 \p@ 前缀为空的标准环境，无节号前缀）
\newtheorem{ithm}{定理}
\newtheorem{iprop}{命题}
\newtheorem{ilem}{引理}
\newtheorem{icor}{推论}
\newtheorem{iconj}{猜想}
\newtheorem{idefn}{定义}
\newtheorem{irem}{注}
\newtheorem{iexam}{例}
\newtheorem{iques}{问题}
\newtheorem{iclaim}{断言}
\newtheorem{inotn}{记号}
\newtheorem{ihyp}{假设}
\newtheorem{ifact}{事实}
\newtheorem{iprob}{问题}
\newenvironment{thm}[1]{\renewcommand{\theithm}{#1}\begin{ithm}}{\end{ithm}}
\newenvironment{prop}[1]{\renewcommand{\theiprop}{#1}\begin{iprop}}{\end{iprop}}
\newenvironment{lem}[1]{\renewcommand{\theilem}{#1}\begin{ilem}}{\end{ilem}}
\newenvironment{cor}[1]{\renewcommand{\theicor}{#1}\begin{icor}}{\end{icor}}
\newenvironment{conj}[1]{\renewcommand{\theiconj}{#1}\begin{iconj}}{\end{iconj}}
\newenvironment{defn}[1]{\renewcommand{\theidefn}{#1}\begin{idefn}}{\end{idefn}}
\newenvironment{rem}[1]{\renewcommand{\theirem}{#1}\begin{irem}}{\end{irem}}
\newenvironment{exam}[1]{\renewcommand{\theiexam}{#1}\begin{iexam}}{\end{iexam}}
\newenvironment{ques}[1]{\renewcommand{\theiques}{#1}\begin{iques}}{\end{iques}}
\newenvironment{claim}[1]{\renewcommand{\theiclaim}{#1}\begin{iclaim}}{\end{iclaim}}
\newenvironment{notn}[1]{\renewcommand{\theinotn}{#1}\begin{inotn}}{\end{inotn}}
\newenvironment{hyp}[1]{\renewcommand{\theihyp}{#1}\begin{ihyp}}{\end{ihyp}}
\newenvironment{fact}[1]{\renewcommand{\theifact}{#1}\begin{ifact}}{\end{ifact}}
\newenvironment{prob}[1]{\renewcommand{\theiprob}{#1}\begin{iprob}}{\end{iprob}}

% 证明环境
\renewenvironment{proof}[1][证明]{\par\medskip\noindent\textit{#1.}\hspace{0.4em}}{\hfill$\square$\par\medskip}

% ---------- 参考文献 ----------
\newenvironment{refs}{\begin{list}{}{\leftmargin=3.8em\labelwidth=3.6em\labelsep=0.4em
  \itemindent=-3.8em\parsep=0.1em\topsep=0.3em\small\raggedright}}{\end{list}}
% 文献条目里作者名与英文标题很长，两端对齐必然溢出，故整段左对齐（不改变内容）。

% ---------- 译注 ----------
\newcommand{\tn}[1]{\footnote{译者注：#1}}
\newcommand{\ednote}[1]{\footnote{编者注：#1}}

% ---------- 数学常用宏（与原文一致） ----------
\newcommand{\C}{\mathbb{C}}
\newcommand{\Q}{\mathbb{Q}}
\newcommand{\Z}{\mathbb{Z}}
\newcommand{\F}{\mathbb{F}}
\newcommand{\R}{\mathbb{R}}
\newcommand{\N}{\mathbb{N}}
\newcommand{\PP}{\mathbb{P}}
\newcommand{\GG}{\mathbb{G}}
\newcommand{\OO}{\mathcal{O}}
\DeclareMathOperator{\Gal}{Gal}
\DeclareMathOperator{\Frobq}{Frob}
\DeclareMathOperator{\Spec}{Spec}
\DeclareMathOperator{\Qcoh}{Qcoh}
\providecommand{\QC}{\Qcoh}% 兼容分片里的简写
\DeclareMathOperator{\Mod}{Mod}
\DeclareMathOperator{\Rep}{Rep}
\DeclareMathOperator{\Ind}{Ind}
\DeclareMathOperator{\End}{End}
\DeclareMathOperator{\Aut}{Aut}
\DeclareMathOperator{\Hom}{Hom}
\DeclareMathOperator{\Ext}{Ext}
\DeclareMathOperator{\Tor}{Tor}
\DeclareMathOperator{\Sym}{Sym}
\DeclareMathOperator{\Stab}{Stab}
\DeclareMathOperator{\Sect}{Sect}
\DeclareMathOperator{\Sing}{Sing}
\DeclareMathOperator{\Gr}{Gr}
\DeclareMathOperator{\Bun}{Bun}
\DeclareMathOperator{\Loc}{Loc}
\DeclareMathOperator{\LocSys}{LocSys}
\DeclareMathOperator{\Perf}{Perf}
\DeclareMathOperator{\IndCoh}{IndCoh}
\DeclareMathOperator{\Coh}{Coh}
\DeclareMathOperator{\QCoh}{QCoh}
\DeclareMathOperator{\Map}{Map}
\DeclareMathOperator{\Spf}{Spf}
\DeclareMathOperator{\Proj}{Proj}
\DeclareMathOperator{\Spa}{Spa}
\DeclareMathOperator{\Sp}{Sp}
\DeclareMathOperator{\HH}{HH}
\DeclareMathOperator{\HC}{HC}
\DeclareMathOperator{\TP}{TP}
\DeclareMathOperator{\TC}{TC}
\DeclareMathOperator{\Frob}{Frob}
\newcommand{\LG}{{}^L G}
\newcommand{\hatG}{\widehat{G}}
\newcommand{\frakg}{\mathfrak{g}}
\newcommand{\frakn}{\mathfrak{n}}
\newcommand{\frakb}{\mathfrak{b}}
\newcommand{\frakt}{\mathfrak{t}}
\newcommand{\frakh}{\mathfrak{h}}
\newcommand{\frakm}{\mathfrak{m}}
\newcommand{\frakp}{\mathfrak{p}}
\newcommand{\frakz}{\mathfrak{z}}
\newcommand{\hG}{\widehat{G}}
\newcommand{\hZ}{\widehat{Z}}
\newcommand{\ol}{\overline}
\newcommand{\wt}{\widetilde}
\newcommand{\wh}{\widehat}
\newcommand{\CC}{\mathcal{C}}
\newcommand{\DD}{\mathcal{D}}
\newcommand{\FF}{\mathcal{F}}
\newcommand{\GGG}{\mathcal{G}}
\newcommand{\LL}{\mathcal{L}}
\newcommand{\MM}{\mathcal{M}}
\newcommand{\NN}{\mathcal{N}}
\newcommand{\TS}{\mathcal{T}}
\newcommand{\XX}{\mathcal{X}}
\newcommand{\ZZ}{\mathcal{Z}}
\newcommand{\isom}{\xrightarrow{\ \sim\ }}
\newcommand{\into}{\hookrightarrow}
\newcommand{\onto}{\twoheadrightarrow}
\newcommand{\ra}{\rightarrow}
\newcommand{\lra}{\longrightarrow}
\newcommand{\dfn}{:=}
\newcommand{\eqdef}{\overset{\mathrm{def}}{=}}
